\section{CAFs as Biomarkers in Breast Cancer}
\label{sec:biomarkers}

The heterogeneity and functional diversity of CAFs suggest their potential as biomarkers for diagnosis, prognosis, and treatment response prediction in breast cancer.

\subsection{Prognostic biomarkers}

Multiple CAF-related features have been associated with clinical outcomes in breast cancer:

\subsubsection{CAF density and distribution}

High overall CAF density---measured by $\alpha$-SMA or FAP staining---is associated with poor prognosis in most breast cancer subtypes \citep{erve2024}. However, the prognostic significance varies by CAF subpopulation:

\textbf{myCAF density:} High myCAF density is associated with aggressive disease features including large tumour size, high grade, and lymph node metastasis \citep{erve2024}. myCAFs contribute to tumour progression through ECM production, mechanical signalling, and immune exclusion, explaining their association with poor outcomes.

\textbf{iCAF density:} The prognostic significance of iCAFs is context-dependent. In some studies, high iCAF density is associated with poor prognosis, reflecting the pro-tumorigenic functions of inflammatory cytokines \citep{erve2024}. In other studies, iCAF density correlates with enhanced immune infiltration and better response to immunotherapy, suggesting a more complex relationship.

\textbf{apCAF density:} Limited data exist on the prognostic significance of apCAFs. Preliminary studies suggest that high apCAF density may be associated with better prognosis in ER+ breast cancer, potentially reflecting enhanced anti-tumour immune responses \citep{erve2024}.

\subsubsection{Spatial organization}

Beyond density, the spatial organization of CAFs provides prognostic information:

\textbf{CAF--tumour distance:} Short distance between CAFs and tumour cells is associated with poor prognosis, reflecting the importance of paracrine signalling in tumour promotion \citep{erve2024}.

\textbf{CAF clustering:} Highly clustered CAF distributions---indicating organized stromal niches---are associated with aggressive disease, while dispersed CAF distributions are associated with better outcomes \citep{erve2024}.

\textbf{Immune--CAF spatial relationships:} The spatial relationship between CAFs and immune cells provides prognostic information. Tumours with CAFs excluded from immune cell-rich areas (immune-excluded phenotype) have worse outcomes than tumours with CAFs integrated within immune infiltrates (immune-infiltrated phenotype) \citep{chhabra2023}.

\subsection{Predictive biomarkers}

CAF features may predict response to specific therapies:

\subsubsection{Immunotherapy response prediction}

The immune-excluded phenotype---characterized by high myCAF density surrounding tumour cell clusters---is associated with poor response to immune checkpoint inhibitors \citep{chhabra2023}. This phenotype can be identified by spatial analysis of $\alpha$-SMA+ myCAFs relative to CD8+ T cells.

Conversely, tumours with high iCAF density and dispersed CAF distributions (immune-inflamed phenotype) show better response to immunotherapy \citep{chhabra2023}. This suggests that CAF subtyping and spatial analysis may help identify patients most likely to benefit from ICI therapy.

\subsubsection{Chemotherapy response prediction}

CAF-related features may predict response to chemotherapy:

\textbf{Stromal gene signatures:} Gene expression signatures derived from CAFs---including ECM components, TGF-$\beta$ targets, and inflammatory cytokines---have been associated with resistance to neoadjuvant chemotherapy in breast cancer \citep{farinha2019}.

\textbf{Collagen organization:} The organization of collagen fibres---measured by second harmonic generation imaging or spatial transcriptomics---is associated with chemotherapy response. Tumours with aligned, dense collagen (reflecting high myCAF activity) show poor response to chemotherapy \citep{provenzano2012}.

\textbf{CAF density:} High overall CAF density is associated with resistance to neoadjuvant chemotherapy, reflecting the physical barrier and paracrine survival signalling provided by CAFs \citep{farinha2019}.

\subsection{Circulating CAF biomarkers}

Circulating CAFs and CAF-derived factors represent a minimally invasive source of biomarkers:

\subsubsection{Circulating fibroblasts}

Circulating fibroblast-like cells---termed circulating fibrocytes---have been detected in the peripheral blood of breast cancer patients \citep{abe2012}. High circulating fibrocyte counts are associated with advanced disease and poor prognosis, potentially reflecting the systemic mobilization of CAF precursors.

However, the detection and characterization of circulating fibrocytes remains technically challenging, and their clinical utility as biomarkers requires further validation.

\subsubsection{Circulating CAF-derived factors}

CAFs secrete multiple factors that can be measured in peripheral blood:

\textbf{CA125 (MUC16):} While traditionally associated with ovarian cancer, CA125 is also secreted by CAFs in breast cancer \citep{seth2021}. Elevated serum CA125 levels are associated with advanced disease and poor prognosis.

\textbf{Tenascin-C:} This ECM glycoprotein, produced by CAFs, is elevated in the serum of breast cancer patients and correlates with disease stage and metastatic burden \citep{seth2021}.

\textbf{Periostin:} This matricellular protein, secreted by myCAFs, is elevated in the serum of breast cancer patients and is associated with poor prognosis \citep{seth2021}.

\subsubsection{Circulating exosomes}

CAFs release exosomes into the bloodstream that carry CAF-specific cargo (miRNAs, proteins) \citep{zhao2023}. These circulating exosomes can be isolated and characterized, providing information about CAF activity without the need for tissue biopsy.

CAF-derived exosomal miRNAs---including miR-21, miR-155, and miR-200 family members---are elevated in the plasma of breast cancer patients and correlate with disease stage and prognosis \citep{zhao2023}.

\subsection{Imaging-based CAF biomarkers}

Non-invasive imaging techniques can capture CAF-related features:

\subsubsection{MRI-based stromal features}

Magnetic resonance imaging (MRI) can assess stromal features that reflect CAF activity:

\textbf{Apparent diffusion coefficient (ADC):} Low ADC values---indicating restricted water diffusion---correspond to dense, collagen-rich stroma produced by myCAFs. Low ADC is associated with poor prognosis and chemoresistance in breast cancer \citep{seth2021}.

\textbf{Dynamic contrast enhancement (DCE):} The pattern of contrast enhancement reflects vascular perfusion, which is influenced by CAF-mediated vascular remodelling. Rapid wash-in and wash-out patterns are associated with high CAF activity and poor prognosis \citep{seth2021}.

\subsubsection{PET-based CAF imaging}

Positron emission tomography (PET) with FAP-targeted tracers ($^{68}$Ga-FAPI) can visualize CAF-rich stromal compartments \citep{kratochwil2019}. High FAPI uptake is associated with dense stroma and may predict resistance to immunotherapy.

\subsection{Multi-modal CAF biomarker panels}

The heterogeneity of CAFs suggests that single biomarkers may be insufficient to capture their full biological complexity. Multi-modal biomarker panels---integrating CAF density, subpopulation composition, spatial organization, and circulating factors---may provide more accurate prognostic and predictive information.

Preliminary studies have shown that combining CAF-related features with tumour-intrinsic features (gene expression, mutation status) improves the prediction of treatment response and clinical outcomes \citep{seth2021}. However, the development and validation of such multi-modal panels requires large, well-annotated clinical cohorts with comprehensive molecular and spatial data.
